\documentclass{beamer}
\usetheme{CambridgeUS}
\definecolor{oxfordblue}{rgb}{0.0, 0.13, 0.28}
\usecolortheme[named=oxfordblue]{structure}
\usepackage{graphicx}
\graphicspath{{./images/}}
\title{Computer Security Hazards}
\author{Anubhav Mehra}
\date{S.S.J University, Almora}
\institute{\textbf{Computer Science}}
\begin{document}
\begin{frame}
	\titlepage
\end{frame}
\begin{frame}{What are Computer Security Hazards}
	\only<1->{
	\begin{block}{Computer Security Hazard}	
		A computer security hazard is really anything on your computer that may damage or steal your data or allow someone else to access your computer, without your knowledge or consent. 
	\end{block}
}
	\begin{itemize}
	
		\item<uncover@2->{There are a lot of different things that can create a computer hazard, like malwares. These include \textbf{\textit{ viruses, worms, ransomware, spyware, and Trojan horses}}.}
		\item<uncover@3->{Misconfiguration of computer products as well as unsafe computing habits also pose risks.}
		\item<uncover@4->{Let's look at these in more detail.}
	\end{itemize}

\end{frame}
\begin{frame}{Malicious Softwares}
	\only<1>{
		Based on Method of Attack
	\begin{itemize}
		\item{A \textbf{worm} is a standalone piece of malicious software that reproduces itself and spreads from computer to computer.}
		\item{A \textbf{virus} is a piece of computer code that inserts itself within the code of another standalone program, then forces that program to take malicious action and spread itself.}
		\item{A \textbf{trojan} is a program that cannot reproduce itself but masquerades as something the user wants and tricks them into activating it so it can do its damage and spread.}
	\end{itemize}
}
\only<2>{
	Based on Purpose/Action
	\begin{itemize}
		\item{A \textbf{spyware} is defined as malware used for the purpose of secretly gathering data on an unsuspecting user.}
		\item{A \textbf{ransomware} is a flavor of malware that encrypts your hard drive's files and demands a payment, usually in Bitcoin, in exchange for the decryption key.}
		\item{A \textbf{rootkit} is, a program or, more often, a collection of software tools that gives a threat actor remote access to and control over a computer or other system.}
	\end{itemize}
}
\end{frame}
\begin{frame}{WannaCry Ransomware Attacks}
	\begin{itemize}
		\only<1-3>{
		\item{WannaCry is a ransomware worm that spread rapidly through across a number of computer networks in May of 2017. After infecting a Windows computers, it encrypts files on the PC's hard drive, making them impossible for users to access, then demands a ransom payment in bitcoin in order to decrypt them.
			}
		}
		\only<2-3>{
		\item{A number of factors made the initial spread of WannaCry particularly noteworthy: it struck a number of important and high-profile systems, including many in Britain's National Health Service;}
		}
		\only<3>{
	\item{The WannaCry ransomware consists of multiple components. It arrives on the infected computer in the form of a dropper, a self-contained program that extracts the other application components embedded within itself.
}
}
\end{itemize}
\end{frame}


\begin{frame}{Misconfiguration Risks}
	Some Misconfiguration Examples are:
	\begin{itemize}
		\only<1->{
		\item{Server misconfigurations like: using \textbf{default} configuration, not disabling \textbf{directory listings}, etc.}
		}
		\only<2->
		{
		\item{Default Operating system configurations like: \textbf{cloud backup}, some \textbf{unsafe applications}, etc.}
		}
		\only<3->
		{
		\item{Software misconfiguraitions like: enabling \textbf{macros is MS Office} , enabling \textbf{javascript in pdfs}, etc.}
		}

	\end{itemize}
	\vspace*{1cm}
	\only<4>{
		How North Korean Hackers Used Word to Hack the Russians:
		\hyperlink{https://www.youtube.com/watch?v=1Fn_mhBzMkQ}{https://www.youtube.com/watch?v=1Fn\_mhBzMkQ}
	}
\end{frame}
\begin{frame}{Unsafe Computing Practices}
	Some Common Unsafe Computing Practices are:
	\begin{itemize}
		\only<1->{
		\item{Using \textbf{weak} passwords or \textbf{same} password for all sites.
		}
	}
		\only<2->{
		\item{Opening Email from \textbf{unverified} party. Going to \textbf{unverified} websites.
		}
	}
		\only<3->{
		\item{Using \textbf{out-dated} technology.}
		}
		\only<4->{
		\item{
				Using \textbf{public} networks like wifi hotspot to do personal transactions.
		}
	}
	\end{itemize}
\end{frame}
\begin{frame}{Basics of Staying Secure}
	\begin{itemize}
		\only<1->{
		\item{Believe that anything online can be \textbf{hacked}.}
		}
		\only<2->{
		\item{Use \textbf{strong} passwords or use \textbf{password generators and managers}.}
		}
		\only<3->{
		\item{Don't open email or files from \textbf{unvarified} sources. Don't visit \textbf{unvarified} sites.}
		}
		\only<4->{
		\item{Keep regular \textbf{offline} backup of your important files and folders. Use encrypted drives if you can.}
		}
		\only<5->{
		\item{Avoid misconfiguraitions in your software-stack. Don't enable \textbf{macros} in MS Office.}
		}
		\only<6->{
		\item{Don't login to any site or transfer any files when using \textbf{public} networks or computers.}
		}
	\end{itemize}
\end{frame}
\begin{frame}{Question and Answers}
	\begin{center}
		\includegraphics[scale=0.2]{thank.jpg}
	\end{center}
	\large Slides for todays presentation can be downloaded from 
	\tiny \textbf{\underline{https://github.com/AnubhavMehraCS/Paper1Presentation/raw/master/ComputerHazard.pdf}}
\end{frame}
\end{document}
